% DeepPass68: detailed differential-geometry derivations.

\begin{derivation}[从\eqnrefs{eq:exterior-derivative-definition-dp68}到\eqnrefs{eq:exterior-calculus-identities-dp68}]
外微分的两个基本性质都来自同一件事：普通偏导对光滑函数可以交换，而基一形式的楔积在交换次序时变号。对 $p$ 阶形式 $\omega$ 再作用一次外微分，
\begin{equation*}
 \dd^2 \omega
 =\frac1{p!}\partial_\rho\partial_\nu\omega_{\mu_1\cdots\mu_p}
 \dd x^\rho\wedge\dd x^\nu\wedge
 \dd x^{\mu_1}\wedge\cdots .
\end{equation*}
只保留 $(\rho,\nu)$ 的反对称部分即可，因为楔积会消去对称部分：
\begin{equation*}
 \partial_\rho\partial_\nu\omega\,
 \dd x^\rho\wedge\dd x^\nu
 =\frac12(\partial_\rho\partial_\nu-\partial_\nu\partial_\rho)\omega\,
 \dd x^\rho\wedge\dd x^\nu=0.
\end{equation*}
所以 $\dd^2=0$。

再令 $\alpha^{(p)}$ 为 $p$ 阶形式。对 $\alpha^{(p)}\wedge\beta$ 的系数求导时，导数落在 $\beta$ 上所产生的新基一形式必须穿过 $p$ 个属于 $\alpha$ 的基一形式，因此得到 $(-1)^p$：
\begin{equation*}
 \dd(\alpha^{(p)}\wedge\beta)
 =(\dd\alpha^{(p)})\wedge\beta
 +(-1)^p\alpha^{(p)}\wedge\dd\beta.
\end{equation*}
两条结果合在一起即为正文\eqnrefs{eq:exterior-calculus-identities-dp68}。
\end{derivation}

\begin{derivation}[从\eqnrefs{eq:connection-oneform-dp11}到\eqnrefs{eq:curvature-components-dp68}]
对表示空间取值的零形式 $\Psi$ 连续作用两次协变微分，
\begin{align*}
 D^2\Psi
 &=\dd(\dd\Psi+\ii g_{\rm G}\mathscr C\Psi)
 +\ii g_{\rm G}\mathscr C\wedge(\dd\Psi+\ii g_{\rm G}\mathscr C\Psi).
\end{align*}
用正文\eqnrefs{eq:exterior-calculus-identities-dp68} 中的 $\dd^2=0$ 和一形式的分次 Leibniz 规则，
\begin{equation*}
 \dd(\mathscr C\Psi)=(\dd\mathscr C)\Psi-\mathscr C\wedge\dd\Psi,
\end{equation*}
于是两个含 $\mathscr C\wedge\dd\Psi$ 的交叉项相消，得到
\begin{align*}
 D^2\Psi
 &=\ii g_{\rm G}(\dd\mathscr C)\Psi
 -g_{\rm G}^2\mathscr C\wedge\mathscr C\,\Psi\\
 &=\ii g_{\rm G}
 \left(\dd\mathscr C+\ii g_{\rm G}\mathscr C\wedge\mathscr C\right)\Psi.
\end{align*}
这就给出正文\eqnrefs{eq:curvature-form-dp11}。

再分别展开 $\dd\mathscr C$ 与 $\mathscr C\wedge\mathscr C$。首先
\begin{align*}
 \dd\mathscr C
 &=(\partial_\mu\mathscr C_\nu)\dd x^\mu\wedge\dd x^\nu\\
 &=\frac12(\partial_\mu\mathscr C_\nu-\partial_\nu\mathscr C_\mu)
 \dd x^\mu\wedge\dd x^\nu.
\end{align*}
其次，矩阵系数不能交换，
\begin{align*}
 \mathscr C\wedge\mathscr C
 &=\mathscr C_\mu\mathscr C_\nu\dd x^\mu\wedge\dd x^\nu\\
 &=\frac12[\mathscr C_\mu,\mathscr C_\nu]\dd x^\mu\wedge\dd x^\nu.
\end{align*}
与 $\mathcal F=\tfrac12\mathcal F_{\mu\nu}\dd x^\mu\wedge\dd x^\nu$ 比较系数，得到正文\eqnrefs{eq:curvature-components-dp68}。
\end{derivation}

\begin{derivation}[从\eqnrefs{eq:covariant-derivative-covariance-dp68}到\eqnrefs{eq:curvature-gauge-covariance-dp68}]
把 $\Psi'=U\Psi$ 代入正文\eqnrefs{eq:covariant-derivative-covariance-dp68}，左边为
\begin{align*}
 D'\Psi'
 &=\dd(U\Psi)+\ii g_{\rm G}\mathscr C'U\Psi\\
 &=(\dd U)\Psi+U\dd\Psi+\ii g_{\rm G}\mathscr C'U\Psi,
\end{align*}
而右边为
\begin{equation*}
 U(D\Psi)=U\dd\Psi+\ii g_{\rm G}U\mathscr C\Psi.
\end{equation*}
消去共同的 $U\dd\Psi$，并利用关系对任意 $\Psi$ 成立，
\begin{equation*}
 \dd U+\ii g_{\rm G}\mathscr C'U=\ii g_{\rm G}U\mathscr C.
\end{equation*}
右乘 $U^{-1}$ 后立即得到正文\eqnrefs{eq:finite-gauge-connection-dp11}。

由刚得到的联络变换式，协变导数满足
\[
D'=UDU^{-1}.
\]
为避免把算符乘法与普通矩阵乘法混淆，让两边作用在任意试探场 $\Psi'$ 上，并写 $\Psi'=U\Psi$。连续作用两次协变导数得到
\begin{align*}
D'^2\Psi'
&=D'\!\left(U D\Psi\right)\\
&=U D(D\Psi)\\
&=U D^2\Psi.
\end{align*}
正文对曲率的定义给出
\[
D^2\Psi=\ii g_{\rm G}\mathcal F\Psi,
\qquad
D'^2\Psi'=\ii g_{\rm G}\mathcal F'\Psi'.
\]
代入 $\Psi'=U\Psi$ 后，两种写法分别为
\begin{align*}
D'^2\Psi'
&=\ii g_{\rm G}\mathcal F' U\Psi,\\
D'^2\Psi'
&=\ii g_{\rm G}U\mathcal F\Psi.
\end{align*}
因为等式对任意 $\Psi$ 成立，且 $U$ 可逆，所以
\[
\mathcal F'U=U\mathcal F,
\qquad
\mathcal F'=U\mathcal F U^{-1},
\]
即正文\eqnrefs{eq:curvature-gauge-covariance-dp68}。
\end{derivation}

\begin{derivation}[从\eqnrefs{eq:covariant-commutator-curvature-dp68}到\eqnrefs{eq:nonabelian-bianchi-dp11}]
Bianchi 恒等式由协变导数对易子的 Jacobi 恒等式直接推出。写出
\begin{equation*}
 [D_\lambda,[D_\mu,D_\nu]]
 +[D_\mu,[D_\nu,D_\lambda]]
 +[D_\nu,[D_\lambda,D_\mu]]=0.
\end{equation*}
由正文\eqnrefs{eq:covariant-commutator-curvature-dp68}，第一项作用在任意 $\Psi$ 上为
\begin{align*}
 [D_\lambda,[D_\mu,D_\nu]]\Psi
 &=\ii g_{\rm G}[D_\lambda,\mathcal F_{\mu\nu}]\Psi\\
 &=\ii g_{\rm G}(D_\lambda\mathcal F_{\mu\nu})\Psi,
\end{align*}
其中
\begin{equation*}
 D_\lambda\mathcal F_{\mu\nu}
 =\partial_\lambda\mathcal F_{\mu\nu}
 +\ii g_{\rm G}[\mathscr C_\lambda,\mathcal F_{\mu\nu}].
\end{equation*}
对另外两项循环置换指标，Jacobi 恒等式便化为
\begin{equation*}
 \ii g_{\rm G}
 (D_\lambda\mathcal F_{\mu\nu}
 +D_\mu\mathcal F_{\nu\lambda}
 +D_\nu\mathcal F_{\lambda\mu})\Psi=0.
\end{equation*}
因为对任意 $\Psi$ 都成立，括号必须为零，即正文\eqnrefs{eq:nonabelian-bianchi-dp11}。
\end{derivation}

\begin{derivation}[从\eqnrefs{eq:parallel-transport-ode-dp68}到\eqnrefs{eq:wilson-endpoint-covariance-dp68}]
记沿路径的矩阵值系数为
\begin{equation*}
 A(s)=\dot x^\mu(s)\mathscr C_\mu[x(s)].
\end{equation*}
把正文\eqnrefs{eq:parallel-transport-ode-dp68} 从 $s_i$ 积到 $s$，得到 Volterra 型积分方程
\begin{equation*}
 \Psi(s)=\Psi(s_i)-\ii g_{\rm G}\int_{s_i}^{s}\dd s_1\,A(s_1)\Psi(s_1).
\end{equation*}
把右边的 $\Psi(s_1)$ 反复代回同一方程，前几阶为
\begin{align*}
 \Psi(s)=\Bigg[I
 &-\ii g_{\rm G}\int_{s_i}^{s}\dd s_1A(s_1)\\
 &+(-\ii g_{\rm G})^2
 \int_{s_i}^{s}\dd s_1\int_{s_i}^{s_1}\dd s_2\,
 A(s_1)A(s_2)\\
 &+(-\ii g_{\rm G})^3
 \int_{s_i}^{s}\dd s_1\int_{s_i}^{s_1}\dd s_2
 \int_{s_i}^{s_2}\dd s_3\,
 A(s_1)A(s_2)A(s_3)+\cdots\Bigg]\Psi(s_i).
\end{align*}
第 $n$ 阶积分域满足 $s_1\ge s_2\ge\cdots\ge s_n$；由于不同 $s$ 处的矩阵一般不交换，较大的路径参数必须排在乘积左侧。用 $\mathcal P$ 记这种顺序，级数重求和为
\begin{equation*}
 \Psi(x_f)=U_C\Psi(x_i),\qquad
 U_C=\mathcal P\exp\!\left[-\ii g_{\rm G}\int_C\mathscr C_\mu\dd x^\mu\right],
\end{equation*}
即正文\eqnrefs{eq:wilson-line-general-dp11}。

现在改变内部空间的局域基底，令 $\Psi'(x)=U(x)\Psi(x)$。同一个输运关系在新基底中满足
\begin{align*}
 \Psi'(x_f)
 &=U(x_f)\Psi(x_f)\\
 &=U(x_f)U_C\Psi(x_i)\\
 &=U(x_f)U_CU^{-1}(x_i)\Psi'(x_i).
\end{align*}
另一方面，新基底中的 Wilson 线 $U_C'$ 按正文定义满足
\begin{equation*}
 \Psi'(x_f)=U_C'\Psi'(x_i).
\end{equation*}
因为该等式对任意初始内部向量 $\Psi'(x_i)$ 都成立，比较线性算符的系数得到
\begin{equation*}
 U_C'=U(x_f)U_CU^{-1}(x_i),
\end{equation*}
即正文\eqnrefs{eq:wilson-endpoint-covariance-dp68}。
\end{derivation}

\begin{derivation}[从\eqnrefs{eq:wilson-line-general-dp11}到\eqnrefs{eq:small-loop-curvature-dp11}]
取左下角为 $x$、两条边向量分别为 $a^\mu$ 与 $b^\nu$ 的小矩形，只保留面积阶 $ab$。四条边的短程 Wilson 线依次为
\begin{align*}
 U_1&=I-\ii g_{\rm G}\mathscr C_\mu(x)a^\mu+\order(a^2),\\
 U_2&=I-\ii g_{\rm G}[\mathscr C_\nu(x)+a^\mu\partial_\mu\mathscr C_\nu(x)]b^\nu
 +\order(b^2,a^2b),\\
 U_3&=I+\ii g_{\rm G}[\mathscr C_\mu(x)+b^\nu\partial_\nu\mathscr C_\mu(x)]a^\mu
 +\order(a^2,ab^2),\\
 U_4&=I+\ii g_{\rm G}\mathscr C_\nu(x)b^\nu+\order(b^2).
\end{align*}
闭合输运为 $U_\square=U_4U_3U_2U_1$。所有一阶项相消；Taylor 项留下
\begin{equation*}
 -\ii g_{\rm G}(\partial_\mu\mathscr C_\nu-\partial_\nu\mathscr C_\mu)a^\mu b^\nu,
\end{equation*}
矩阵乘法的二阶交叉项留下
\begin{equation*}
 -g_{\rm G}^2[\mathscr C_\mu,\mathscr C_\nu]a^\mu b^\nu.
\end{equation*}
合并并使用正文\eqnrefs{eq:curvature-components-dp68}，便得到正文\eqnrefs{eq:small-loop-curvature-dp11}。
\end{derivation}

\begin{derivation}[从\eqnrefs{eq:pontryagin-form-definition-dp68}到\eqnrefs{eq:pontryagin-number-dp11}]
以下计算先核对两种写法的归一化；Pontryagin 数的整数性来自整体拓扑条件，而不是这一局部代数恒等式。写
\begin{equation*}
 \mathcal F=\frac12\mathcal F_{\mu\nu}^aT^a
 \dd x^\mu\wedge\dd x^\nu,
 \qquad
 \Tr(T^aT^b)=\frac12\delta^{ab}.
\end{equation*}
于是
\begin{align*}
 \Tr(\mathcal F\wedge\mathcal F)
 &=\frac18\mathcal F_{\mu\nu}^a\mathcal F_{\rho\sigma}^a
 \epsilon^{\mu\nu\rho\sigma}\dd^4 x.
\end{align*}
再用对偶场强定义
\begin{equation*}
 {}^\star\!\mathcal F^{a\mu\nu}
 =\frac12\epsilon^{\mu\nu\rho\sigma}\mathcal F_{\rho\sigma}^a,
\end{equation*}
可写成
\begin{equation*}
 \Tr(\mathcal F\wedge\mathcal F)
 =\frac14\mathcal F_{\mu\nu}^a
 {}^\star\!\mathcal F^{a\mu\nu}\dd^4 x.
\end{equation*}
代回正文\eqnrefs{eq:pontryagin-form-definition-dp68}，得到\eqnrefs{eq:pontryagin-number-dp11} 中的分量归一化 $g_{\rm G}^2/(32\pi^2)$。其中 $Q_{\rm top}\in\mathbb Z$ 的部分依赖正文已经说明的边界条件与整体拓扑定理，不属于这个代数推导的结论。
\end{derivation}
